\section{Determination of the spectrum}\label{secsig}
%%%%%%%%%%%%%%%%%%%%%%%%%%%%%%%%%%%%%%%%%%%%%%

In this section we aim at the determination of the spectrum $\sigma(M)$.
We start with the case where $M$ is a Hadamard manifold with sectional curvature $K_M<0$
which is asymptotically harmonic about a point $\xi\in M_\infty$
with mean curvature of horospheres equal to $h>0$.
Recall from the first part of the proof of Theorem~\ref{onlyacn}
that the case left is $N=\Gamma\backslash M$,
where the fibers of the associated Riemannian submersion $\pi=b\colon N\to\R$ are of infinite volume,
where $b$ is the push-down to $N$ of some Busemann function associated to $\xi$.
We subsume the case $N=M$ here by including $\Gamma=\{\id\}$.
From Observation~\ref{chee}, we know that $\sigma(N)\subseteq\bigl[h^2/4,\infty\bigr)$.

For any smooth function $v$ on $\R$, we have
\begin{align*}
	\Delta b^*v = b^*\bigl(-v'' - hv'\bigr).
\end{align*}
Now $\vf=\vf(t)=\exp(ht/2)$ satisfies $\vf'=h\vf/2$ and $\vf''=h^2\vf/4$.
With $\mu=\vf^2{\rm d}t$, the unitary transformation
\begin{align*}
	\Phi \colon \ L^2(\R,\mu) \to L^2(\R), \qquad \Phi(v) = \vf v,
\end{align*}
yields
\begin{align*}
	(\Phi(v))'' = (\vf v)'' = \vf\biggl(\frac{h^2}4v + hv' + v''\biggr),
\end{align*}
and hence $\Phi$ intertwines the diffusion operator $Lv=-(v''+hv')$ in $L^2(\R,\mu)$
with the Schr\"odinger operator $Sw=-w''+h^2w/4$ in $L^2(\R)$.
Therefore, the spectra of $L$ and $S$ coincide,
where we recall that $H_{\ac}(S)=L^2(\R)$ with $\sigma(S)=\bigl[h^2/4,\infty\bigr)$.

We want to show that $\sigma(S)\subseteq\sigma(N)$.
In contrast to the case where the volume of the fibers of $b$ are of finite volume,
we cannot simply use pull backs $b^*v$ since we would lose square-integrability that way.
What we will do is to carefully cut off pull-backs.

\begin{Lemma}
Suppose that $-1\le K_M\le-a^2<0$ and that $\|\nabla R_M\|_\infty<\infty$.
Let $\chi_0\colon b^{-1}(0)\to\R$ be a non-zero $C^2$ function with compact support.
Extend $\chi_0$ to $N$ by letting $\chi=\chi_t$ be equal to~$F_{-t}^*\chi_0$ on $b^{-1}(t)$,
where $(F_t)_{t\in\R}$ denotes the flow of $X=\nabla b$.
Then $\chi$ is $C^2$ and $\nabla\chi$ and $\Delta\chi$ tend to zero uniformly as $t\to\infty$.
\end{Lemma}

\begin{proof}
Since $\chi$ is constant along the flow lines of $X$,
the gradient of $\chi$ is tangential to the fibers of $b^{-1}(t)$ of $b$.
By the upper curvature bound $-a^2$, unstable Jacobi field $J$ grow at least exponentially,
\begin{align*}
	|J(t)| \ge {\rm e}^{at} |J(0)| \qquad\text{for all} \quad t>0.
\end{align*}
Therefore,
\begin{align*}%\label{enab}
	|\nabla\chi_t| \le {\rm e}^{-at} |\nabla\chi_0| \qquad\text{for all}\quad t>0,
\end{align*}
which implies the first assertion.

For the estimate of the Laplacian, we refer to~\cite[Section 6]{BP21b}.
The situation there is that of a Hadamard manifold $X$ with pinched negative sectional curvature
and uniformly bounded covariant derivative of its curvature tensor,
a convex domain $C$ in $X$,
and a function obtained by extending a given function on $\partial C$
to $X\setminus C$ constantly along minimizing geodesics to $C$.
In our situation, $X$ corresponds to $M$, $C$ to the horoball $b^{-1}((-\infty,0])$,
and the given function to $\chi_0$.
A short look at the arguments in~\cite[Section 6]{BP21b} shows
that they also apply in the corresponding situation in $N$.
The estimate obtained in~\cite{BP21b} is that
\begin{align*}
	|\Delta\chi| \le c_0{\rm e}^{-at} \qquad\text{over} \quad b^{-1}(t)\quad \text{for all} \quad t>0,
\end{align*}
where $c_0$ is a constant, which depends on bounds for $K_M$, $\nabla R_M$, $\nabla\chi_0$, and $\nabla^2\chi_0$;
compare with~\cite[formula~(6.4) and end of Section~6]{BP21b}.
\end{proof}

\begin{proof}[End of proof of Theorem~\ref{onlyacn}]
Let $\lambda\in\sigma(S)$ and $\ve>0$.
Choose a non-vanishing $w\in C^\infty_{\rm c}(\R)$ such that $\|(S-\lambda)w\|_2\le\ve\|w\|_2$.
By shifting $w$ along $\R$ if necessary, we can assume that $\supp w\subseteq[t,\infty)$,
where $t>0$ is such that $|\Delta\chi| < \ve$ on $b^{-1}([t,+\infty))$.
Note that shifting $w$ is legitimate since $Sw$ is only shifted accordingly.
Then $v=w/\vf$ also has support in $[t,\infty)$ and satisfies $\|(L-\lambda)v\|_2\le\ve\|v\|_2$,
where now the $L^2$-norm is taken with respect to $\mu=\vf^2{\rm d}t$.
Then
\begin{align*}
	\Delta(\chi b^*v)
	= (\Delta\chi)b^*v - 2\langle\nabla\chi,\nabla b^*v\rangle + \chi\Delta b^*v
	= (\Delta\chi)b^*v + \chi b^*Lv
\end{align*}
since $\nabla\chi$ is perpendicular to $X$ and $\nabla b^*v$ is collinear with $X$
and since $b^*$ intertwines $\Delta$ with $L$.

Since $b$ is a Riemannian submersion and the flow maps $F_s$ of $X$ induces
a diffeomorphism of~$b^{-1}(0)$ with $b^{-1}(s)$,
for all $s\in\R$ and with respective Jacobian $\vf^2(s)={\rm e}^{hs}$.
Thus we identify~$x\in N$ with $(y,s)\in b^{-1}(0)\times\R$, where $b(x)=s$ and $F_s(y)=x$,
and get
\begin{align*}%\label{n1}
	\|\chi b^*v\|_2^2
	&= \int \chi^2(x)(b^*v(x))^2{\rm d}x \\
	&= \iint_{b^{-1}(0)\times\R}\chi_0^2(y)v^2(s)\vf^2(s){\rm d}y{\rm d}s \\
	&= \|\chi_0\|_2^2 \int_\R v^2(s)\vf^2(s){\rm d}s
	= \|\chi_0\|_2^2 \|v\|_2^2,
\end{align*}
where $\|\chi_0\|_2$ denotes the $L^2$-norm of $\chi_0$ in $L^2\bigl(b^{-1}(0)\bigr)$.
Likewise, by the choice of $t$,
\begin{align*}%\label{n2}
	\|(\Delta\chi)b^*v\|_2^2
	\le \ve^2|\supp\chi_0| \|v\|_2^2.
\end{align*}
and hence
\begin{align*}%\label{n3}
	\|(\Delta\chi)b^*v\|_2^2
	\le \ve^2\frac{|\supp\chi_0|}{ \|\chi_0\|_2^2}\|\chi b^*v\|_2^2
	= \ve^2c_1^2\|\chi b^*v\|_2^2,
\end{align*}
where $c_1$ depends on the choice of $\chi_0$.
Hence we obtain
\begin{align*}%\label{n4}
	\|(\Delta-\lambda)(\chi b^*v)\|_2^2
	&\le 2\|\chi b^*(L-\lambda)v\|_2^2 + 2\ve^2c_1^2\|\chi b^*v\|_2^2 \\
	&= 2\|\chi_0\|_2^2 \|(L-\lambda)v\|_2^2 + 2\ve^2c_1^2\|\chi b^*v\|_2^2 \\
	&= \ve^2c_2^2\|\chi b^*v\|_2^2,
\end{align*}
where $c_2$ depends on the choice of $\chi_0$.
This implies that $\lambda\in\sigma(N)$.
\end{proof}
